\documentclass[11pt]{article}
\usepackage[letterpaper,margin=1in,headheight=14pt]{geometry}
\usepackage[T1]{fontenc}
\usepackage{lmodern}
\usepackage{amsmath,amssymb,amsthm,mathtools}
\usepackage{microtype,booktabs,array,xcolor,fancyhdr,xurl}
\usepackage[colorlinks=true,linkcolor=blue!45!black,citecolor=blue!45!black,urlcolor=blue!45!black]{hyperref}
\hypersetup{pdftitle={Gaussian variance monotonicity for de-regularized MMD},pdfauthor={Siavash Sadeghi (contributor and submitter); AI-generated research note},pdfsubject={Proposed affirmative resolution of AIM Problem 560, commit aa776a01d7d48a79f93251af11fde9454b0aea95}}
\newtheorem{theorem}{Theorem}[section]
\newtheorem{lemma}[theorem]{Lemma}
\newtheorem{proposition}[theorem]{Proposition}
\newtheorem{corollary}[theorem]{Corollary}
\theoremstyle{remark}
\newtheorem{remark}[theorem]{Remark}
\newcommand{\R}{\mathbb R}
\newcommand{\E}{\mathbb E}
\newcommand{\Ical}{\mathcal I}
\newcommand{\ip}[2]{\left\langle #1,#2\right\rangle}
\newcommand{\dd}{\,\mathrm d}
\newcommand{\Var}{\operatorname{Var}}
\newcommand{\diag}{\operatorname{diag}}
\pagestyle{fancy}
\fancyhf{}
\fancyhead[L]{\small Gaussian variance monotonicity for DrMMD}
\fancyhead[R]{\small AIM 560 \;|\; proposed resolution}
\fancyfoot[C]{\thepage}
\renewcommand{\headrulewidth}{0.3pt}
\setlength{\parindent}{0pt}
\setlength{\parskip}{5.5pt}
\setlength{\emergencystretch}{2em}
\allowdisplaybreaks[2]
\begin{document}
\thispagestyle{empty}
{\small\sffamily\bfseries RESEARCH NOTE \quad AIM 560 \quad 1 OCTOBER 2026}
\vspace{8pt}
\begin{center}
{\LARGE\bfseries Gaussian variance monotonicity\par}
\vspace{4pt}
{\LARGE\bfseries for de-regularized MMD\par}
\vspace{10pt}
{\large A finite-projection proof for every positive regularization parameter}
\end{center}
\vspace{4pt}
\begin{center}
Contributor and submitter: \textbf{Siavash Sadeghi}
\end{center}
\textbf{Status.} Complete proposed affirmative resolution; independent mathematical review pending.
This AI-generated note is not a publication, an independent audit, or a proof-assistant certificate.
The target is AIM Problem 560 at repository commit \path{aa776a01d7d48a79f93251af11fde9454b0aea95}, inspected on 1 October 2026.

\begin{abstract}
For a centered one-dimensional Gaussian target and the Gaussian kernel, we prove the strict sign inequality governing the variance in the Gaussian moment approximation to de-regularized maximum mean discrepancy (DrMMD) flow. The result holds for every target variance $s>0$, every finite regularization parameter $\lambda>0$, and every current variance $0<v<s$, removing the target statement's restriction $s>2$. The proof reduces the Gaussian expectations to a weighted power-series curve. Every finite initial block of its coordinates has a strictly negative directional inner product toward the target. A truncated negative-binomial mean bound proves this fact for positive curve parameters; a polynomial boundary identity proves it for zero and negative parameters. Summation by parts then handles all nonnegative decreasing spectral weights, including the DrMMD resolvent. The Gaussian moment variance consequently increases to its target. The analytic argument does not rely on the accompanying computations.
\end{abstract}

\section{The exact target and the result}
Fix $s>0$ and write $\pi_s=N(0,s)$. On $L^2(\pi_s)$ let
\begin{equation}\label{eq:T}
 (T_s f)(y)=\int_\R e^{-(x-y)^2/2}f(x)\dd\pi_s(x).
\end{equation}
For $0<v<s$ set $\nu_v=N(0,v)$,
$f_v=\mathrm d\nu_v/\mathrm d\pi_s$, and $r_v=f_v-1$.
For $\lambda>0$ define the normalized witness
\begin{equation}\label{eq:witness}
 h_{v,\lambda}=T_s(T_s+\lambda I)^{-1}r_v.
\end{equation}
AIM Problem 560 asks whether the quantity
\begin{equation}\label{eq:target}
 \Ical(s,v,\lambda)=\int_\R y\,h_{v,\lambda}'(y)\dd\nu_v(y)
\end{equation}
is strictly negative for $s>2$, $\lambda>0$, and $0<v<s$ \cite{aim}.
The question comes from the finite-regularization conjecture in Appendix C of Chen et al.\ \cite{chen}; that source establishes the two limiting regimes.

\begin{theorem}\label{thm:main}
For every $s>0$, $\lambda>0$, and $0<v<s$, the witness in \eqref{eq:witness} satisfies
\[
 \boxed{\Ical(s,v,\lambda)<0.}
\]
Consequently the full parameter range requested in AIM Problem 560 satisfies the conjectured inequality.
\end{theorem}
The proof concerns the \emph{Gaussian moment approximation}. It does not assume or assert that the unrestricted DrMMD transport flow preserves Gaussian distributions.

\section{A rigorous reduction to a scalar spectral series}
\subsection{Regularity and an integration-by-parts identity}
The density ratio is
\[
 f_v(y)=\sqrt{\frac{s}{v}}\exp\!\left[-\frac12\left(\frac1v-\frac1s\right)y^2\right].
\]
The map $v\mapsto f_v$ is smooth as a map from $(0,2s)$ into $L^2(\pi_s)$. For a compact interval $[a,b]\subset(0,2s)$ and every nonnegative integer $k$, the squared variance derivatives have the uniform integrable envelope
\[
 |\partial_v^k f_v(y)|^2\frac{\mathrm d\pi_s}{\mathrm dy}(y)
 \le C_k(1+|y|^{4k})\exp\!\left[-\left(\frac1b-\frac1{2s}\right)y^2\right].
\]
The exponent coefficient is positive because $b<2s$. This justifies differentiation in $L^2(\pi_s)$ even when $v>s$.
In particular,
\begin{equation}\label{eq:fvprime}
 f_v'(y)=\frac{y^2-v}{2v^2}f_v(y).
\end{equation}
Here and below a prime on $f_v$ means differentiation in the variance, while a prime on $h$ means differentiation in its spatial argument.

The operator $T_s$ is bounded, positive and self-adjoint. The explicit diagonalization below also verifies positivity. Put $g=(T_s+\lambda I)^{-1}r_v$. Then $g\in L^2(\pi_s)\subset L^1(\pi_s)$. Both $k(x,y)$ and $\partial_y k(x,y)=(x-y)k(x,y)$ are uniformly bounded, so $h=T_sg$ and its spatial derivative are bounded and continuous. Gaussian integration by parts therefore gives
\begin{align}
 \Ical(s,v,\lambda)
 &=\int_\R \left(\frac{y^2}{v}-1\right)h_{v,\lambda}(y)\dd\nu_v(y)\notag\\
 &=2v\,\ip{h_{v,\lambda}}{f_v'}_{L^2(\pi_s)}.\label{eq:ibp}
\end{align}
This avoids any need to differentiate an infinite witness expansion pointwise.

\subsection{Gaussian eigenfunctions, including degree zero}
Introduce
\begin{equation}\label{eq:parameters}
 q=\frac{\sqrt{1+4s}-1}{\sqrt{1+4s}+1}\in(0,1),\qquad
 s=\frac{q}{(1-q)^2},\qquad
 \eta=\frac{1-q}{2},\quad
 \gamma^2=\frac{1-q^2}{2q},\quad
 \beta=\sqrt{\frac{1+q}{1-q}}.
\end{equation}
Let $H_j$ be the physicists' Hermite polynomial, characterized by
$e^{2xt-t^2}=\sum_{j\ge0}H_j(x)t^j/j!$.
Define
\begin{equation}\label{eq:eigenfunctions}
 e_j(x)=\sqrt{\frac{\beta}{2^j j!}}\,e^{-\eta x^2}H_j(\gamma x),
 \qquad \rho_j=(1-q)q^j,\qquad j=0,1,2,\ldots.
\end{equation}
These agree, after reparameterization, with the Gaussian spectral formulas underlying Appendix C of \cite{chen}. For completeness, their normalization and eigenvalue equation are verified here.

Since
\[
 2\eta+\frac1{2s}=\gamma^2,
 \qquad \sqrt{2s}\,\gamma=\beta,
\]
Hermite orthogonality under $e^{-u^2}\dd u$ gives
$\ip{e_i}{e_j}_{L^2(\pi_s)}=\delta_{ij}$.
The same change of variable and multiplication by the Gaussian factor transfers completeness from the Hermite basis, so $(e_j)_{j\ge0}$ is an orthonormal basis of $L^2(\pi_s)$.

A direct Gaussian integral, using
$1+2s(\eta+1/2)=(1-q)^{-2}$, yields
\begin{align*}
 &T_s\!\left(e^{-\eta x^2}e^{2\gamma xt-t^2}\right)(y)\\
 &\hspace{12pt}=(1-q)e^{-\eta y^2}e^{2\gamma y(qt)-(qt)^2}.
\end{align*}
Comparing Taylor coefficients in $t$ proves $T_s e_j=\rho_j e_j$.
In particular $T_s$ is positive. Notice that $e_0(x)=\sqrt\beta e^{-\eta x^2}$ is \emph{not} the constant function. The degree-zero eigenmode must not be omitted.

\subsection{Explicit Gaussian averages}
Set
\begin{equation}\label{eq:z}
 z=z(v)=\frac{(1-q)v/q-1}{1+(1-q)v}.
\end{equation}
Then
\begin{equation}\label{eq:zmap}
 -1<z<q,\quad z(s)=q,\quad
 z'(v)=\frac{1-q^2}{q(1+(1-q)v)^2}>0,
 \quad 1+(1-q)v=\frac{1+q}{1-qz}.
\end{equation}
Define
\begin{equation}\label{eq:phi}
 c_n=\frac{\binom{2n}{n}}{4^n},\qquad
 \phi_n(t)=\sqrt{c_n}\sqrt{1-qt}\,t^n,\qquad
 K=(1-q^2)^{-1/4}.
\end{equation}
For $m_j(v)=\int e_j\dd\nu_v$, symmetry gives $m_{2n+1}(v)=0$, while
\begin{equation}\label{eq:averages}
 m_{2n}(v)=K\phi_n(z),\qquad m_{2n}(s)=K\phi_n(q).
\end{equation}
Indeed, the generating integral is
\[
 \E_{Y\sim N(0,v)}\!\left[e^{-\eta Y^2}e^{2\gamma Yt-t^2}\right]
 =\frac{e^{z t^2}}{\sqrt{1+(1-q)v}}.
\]
Coefficient extraction at degree $2n$, followed by \eqref{eq:zmap}, gives \eqref{eq:averages}.

The even-mode resolvent weights are
\begin{equation}\label{eq:weights}
 w_n=\frac{\rho_{2n}}{\rho_{2n}+\lambda}
 =\frac{(1-q)q^{2n}}{(1-q)q^{2n}+\lambda}.
\end{equation}
Thus $w_n>0$, $w_{n+1}<w_n$, and $w_n\to0$.
Using \eqref{eq:ibp} and Parseval's identity gives the precise reduction
\begin{equation}\label{eq:reduction}
 \boxed{\Ical(s,v,\lambda)=2vK^2z'(v)
 \sum_{n=0}^\infty w_n\bigl(\phi_n(z)-\phi_n(q)\bigr)\phi_n'(z).}
\end{equation}
The prime on $\phi_n$ denotes its real scalar argument.
The multiplier outside the series is strictly positive.

\section{The finite-prefix sign lemma}
All functions in this section have a fixed $q\in(0,1)$.
Put
\begin{equation}\label{eq:SN}
 S_N(q,z)=\sum_{n=0}^N\bigl(\phi_n(z)-\phi_n(q)\bigr)\phi_n'(z),
 \quad
 P_N(u)=\sum_{n=0}^N c_nu^n,\quad C_N=(2N+1)c_N.
\end{equation}
The decisive point is that every $S_N$ is negative, even though individual summands can be positive.

\begin{lemma}[Finite-prefix sign]\label{lem:prefix}
For every integer $N\ge0$, $0<q<1$, and $-1<z<q$,
\[
 S_N(q,z)<0.
\]
\end{lemma}

\subsection{Two elementary coefficient identities}
The recurrence $c_{n+1}/c_n=(2n+1)/(2n+2)$ gives
\begin{equation}\label{eq:poly}
 2(1-u)P_N'(u)-P_N(u)=-C_Nu^N.
\end{equation}
Also,
\begin{equation}\label{eq:Cbound}
 C_N^2\le 2N+1.
\end{equation}
To see this, start with $c_0=1$ and use
\[
 (2N+1)c_N^2=(2N-1)c_{N-1}^2\left(1-\frac1{4N^2}\right)
 \quad(N\ge1).
\]
Hence $(2N+1)c_N^2\le1$, which is equivalent to \eqref{eq:Cbound}.
For $0<u<1$,
\begin{equation}\label{eq:sos}
 \sum_{k=0}^{2N}u^k-(2N+1)u^N
 =\sum_{k=0}^{N-1}u^k(1-u^{N-k})^2\ge0.
\end{equation}
The empty sum is zero when $N=0$.
Combining \eqref{eq:Cbound} and \eqref{eq:sos}, with $u=q^2$, yields
\begin{equation}\label{eq:boundarybound}
 C_N^2(1-q^2)q^{2N}
 \le (2N+1)(1-q^2)q^{2N}
 \le1-q^{4N+2}<1.
\end{equation}

\subsection{Positive curve parameter: a truncated mean inequality}
Suppose $0<z<q$.
Define the finite kernel
\[
 G_N(x,y)=\sum_{n=0}^N\phi_n(x)\phi_n(y)
 =\sqrt{(1-qx)(1-qy)}\,P_N(xy).
\]
We first show that
\begin{equation}\label{eq:mixedpositive}
 \partial_x\partial_y G_N(x,y)>0
 \quad\text{whenever }0<x,y\le q.
\end{equation}
For $u=xy$, define a probability mass function on $\{0,\ldots,N\}$ by
\[
 p_n=\frac{c_nu^n}{P_N(u)},\qquad
 \mu=\sum_{n=0}^N n p_n,\qquad
 V=\sum_{n=0}^N(n-\mu)^2p_n.
\]
Equation \eqref{eq:poly} implies
\begin{equation}\label{eq:mean}
 \mu=\frac{uP_N'(u)}{P_N(u)}
 =\frac{u}{2(1-u)}\left(1-\frac{C_Nu^N}{P_N(u)}\right)
 <\frac{u}{2(1-u)}.
\end{equation}
This is also the mean bound obtained by truncating a negative-binomial distribution with shape $1/2$.
Set
\[
 a_x=\frac{qx}{2(1-qx)},\qquad a_y=\frac{qy}{2(1-qy)}.
\]
Because $xy\le qx$ and $xy\le qy$, and $t/(2(1-t))$ increases on $(0,1)$, equation \eqref{eq:mean} gives $\mu<a_x$ and $\mu<a_y$.
For $x>0$,
$\phi_n'(x)=\phi_n(x)(n-a_x)/x$.
Consequently
\begin{align}
 \partial_x\partial_yG_N(x,y)
 &=\frac{\sqrt{(1-qx)(1-qy)}P_N(xy)}{xy}
 \sum_{n=0}^N p_n(n-a_x)(n-a_y)\notag\\
 &=\frac{\sqrt{(1-qx)(1-qy)}P_N(xy)}{xy}
 \left[V+(\mu-a_x)(\mu-a_y)\right]>0.\label{eq:covariance}
\end{align}
This proves \eqref{eq:mixedpositive}, including $N=0$.
Finally,
\begin{equation}\label{eq:positiveprefix}
 S_N(q,z)=-\int_z^q\partial_x\partial_yG_N(z,t)\dd t<0.
\end{equation}

\subsection{Zero and negative curve parameter: a boundary identity}
The derivative of the cross kernel at the target simplifies exactly:
\begin{align}
 \frac{\mathrm d}{\mathrm dz}G_N(z,q)
 &=\frac{q\sqrt{1-q^2}}{2\sqrt{1-qz}}
 \left[2(1-qz)P_N'(qz)-P_N(qz)\right]\notag\\
 &=-\frac{q\sqrt{1-q^2}}{2\sqrt{1-qz}}C_N(qz)^N.\label{eq:crossderivative}
\end{align}
Using half the derivative of $G_N(z,z)$ for the self term gives
\begin{equation}\label{eq:prefixformula}
 S_N(q,z)=(1-qz)zP_N'(z^2)-\frac q2P_N(z^2)
 +\frac{q\sqrt{1-q^2}}{2\sqrt{1-qz}}C_N(qz)^N.
\end{equation}
At $z=0$, this reads
\[
 S_0(q,0)=-\frac q2\bigl(1-\sqrt{1-q^2}\bigr)<0,
 \qquad S_N(q,0)=-\frac q2<0\quad(N\ge1).
\]
Now let $z=-x$ with $0<x<1$.
Since $P_N(x^2)\ge1$ and $P_N'(x^2)\ge0$, the first two terms in \eqref{eq:prefixformula} satisfy
\[
 -x(1+qx)P_N'(x^2)-\frac q2P_N(x^2)\le-\frac q2.
\]
The absolute value of the remaining term is
\[
 \frac q2\left(C_N\sqrt{1-q^2}\,q^N\right)
 \frac{x^N}{\sqrt{1+qx}}<\frac q2
\]
by \eqref{eq:boundarybound}.
It follows that $S_N(q,-x)<0$.
Together with \eqref{eq:positiveprefix} and the $z=0$ calculation, this proves Lemma~\ref{lem:prefix}.\hfill$\square$

\begin{remark}[Why a termwise argument would fail]
Take $q=9/10$, $z=3/5$, and $n=1$.
Then $(1-qz)z^2>(1-q^2)q^2$ and $1-3qz/2>0$, so
$(\phi_1(z)-\phi_1(q))\phi_1'(z)>0$.
The proved property is negativity of every \emph{initial block}, not negativity of every mode separately.
\end{remark}

\section{Every decreasing spectral filter has the sign}
Let
$a_n=(\phi_n(z)-\phi_n(q))\phi_n'(z)$.
The series $\sum_n|a_n|$ converges for $-1<z<q<1$.
For $z\ne0$, this follows from a bound by a constant times
\[
 (n+1)\left(|z|^{2n-1}+q^n|z|^{n-1}\right),\qquad n\ge1;
\]
for $z=0$, all sufficiently high terms vanish.
Thus the prefix sums $S_N$ are bounded and converge.
Using $P_\infty(u)=(1-u)^{-1/2}$ gives
\begin{equation}\label{eq:Sinfty}
 S_\infty(q,z)=\frac{z-q}{2(1-z^2)^{3/2}}<0.
\end{equation}
The cross kernel $G_\infty(z,q)=\sqrt{1-q^2}$ is constant in $z$, which also explains \eqref{eq:Sinfty}.

\begin{proposition}[Decreasing-weight sign]\label{prop:weights}
Suppose $(b_n)_{n\ge0}$ is bounded, nonnegative and nonincreasing, and $b_0>0$.
Then for every $0<q<1$ and $-1<z<q$,
\[
 \sum_{n=0}^\infty b_n(\phi_n(z)-\phi_n(q))\phi_n'(z)<0.
\]
\end{proposition}
\begin{proof}
Let $b_\infty=\lim_n b_n\ge0$.
Finite summation by parts gives
\[
 \sum_{n=0}^M b_na_n
 =\sum_{n=0}^{M-1}(b_n-b_{n+1})S_n+b_MS_M.
\]
Since $S_n$ is bounded and $\sum_n(b_n-b_{n+1})=b_0-b_\infty$, passage to the limit is justified and yields
\begin{equation}\label{eq:abel}
 \sum_{n=0}^\infty b_na_n
 =\sum_{n=0}^\infty(b_n-b_{n+1})S_n+b_\infty S_\infty.
\end{equation}
All coefficients on the right are nonnegative. Their sum is $b_0>0$, and every $S_n$ and $S_\infty$ is strictly negative by Lemma~\ref{lem:prefix} and \eqref{eq:Sinfty}. Therefore the right side is strictly negative.
\end{proof}

\begin{proof}[Proof of Theorem~\ref{thm:main}]
The DrMMD weights \eqref{eq:weights} satisfy Proposition~\ref{prop:weights}.
Equation \eqref{eq:reduction} multiplies their strictly negative series by $2vK^2z'(v)>0$. This proves \eqref{eq:target} is strictly negative for every $s>0$, $\lambda>0$, and $0<v<s$.
\end{proof}

\section{A quantitative bound and the Gaussian moment evolution}
\subsection{A strictly negative explicit upper bound}
For the DrMMD weights, the $N=0$ term alone on the right side of \eqref{eq:abel} provides the bound
\begin{equation}\label{eq:quantbound}
 \Ical(s,v,\lambda)\le
 -\frac{v q z'(v)}{\sqrt{1-q^2}}(w_0-w_1)
 \left(1-\sqrt{\frac{1-q^2}{1-qz}}\right)<0,
\end{equation}
where
\[
 w_0-w_1=
 \frac{\lambda(1-q)(1-q^2)}{(1-q+\lambda)((1-q)q^2+\lambda)}>0.
\]
The other factors are positive because $v>0$, $z'(v)>0$, and $z<q$.
The bound is not intended to be sharp; its purpose is to exhibit a nonzero sign margin in closed form.

\subsection{Monotone convergence in the stated scalar model}
\begin{corollary}\label{cor:flow}
Fix $s>0$ and $\lambda>0$. For every $v_0\in(0,s)$ the scalar initial-value problem
\[
 \dot v=-2(1+\lambda)\Ical(s,v,\lambda),\qquad v(0)=v_0
\]
has a unique global forward solution, for $t\ge0$, with $v_0\le v(t)<s$ for finite $t$, $\dot v(t)>0$, and $v(t)\to s$ as $t\to\infty$.
\end{corollary}
\begin{proof}
By \eqref{eq:ibp}, the right-hand side is smooth on $(0,2s)$ and vanishes at $v=s$.
Theorem~\ref{thm:main} makes it positive for $v<s$.
Local uniqueness prevents a solution starting below $s$ from reaching or crossing the equilibrium at a finite time. The trajectory therefore stays in the compact interval $[v_0,s]$, where the smooth vector field is bounded, and continues globally.
Its increasing limit $L$ is at most $s$. If $L<s$, continuity would give a strictly positive lower bound on the speed near $L$, contradicting convergence to $L$. Hence $L=s$.
\end{proof}
For completeness, writing $A_\lambda=T_s(T_s+\lambda I)^{-1}$, the derivative at equilibrium is
\[
 \left.\frac{\mathrm d}{\mathrm dv}\bigl[-2(1+\lambda)\Ical(s,v,\lambda)\bigr]\right|_{v=s}
 =-4s(1+\lambda)\ip{A_\lambda f_s'}{f_s'}_{L^2(\pi_s)}<0.
\]
The inequality follows because every eigenvalue of $A_\lambda$ is positive and $f_s'\ne0$.
Thus the scalar equilibrium is also locally exponentially attracting.
No corresponding claim is made here for arbitrary distributions, nonzero means, higher dimensions, discrete particle updates, or variances initially above $s$.

\section{Verification, reproducibility and research status}
The proof is analytic: no finite grid or numerical approximation is used to infer the universally quantified theorem.
The accompanying standard-library program \texttt{verify\_exact.py} passed \textbf{2,237 exact supporting checks}, including the coefficient recurrence, the polynomial differential identity, a sum-of-squares identity, representative mixed-derivative identities, and rigorous enclosures for seven infinite spectral sums.
These are supporting checks, not formal verification of the theorem.

The separate program \texttt{verify\_numerics.py} compares the spectral expression against a direct Gaussian-kernel Nystr\"om discretization of \eqref{eq:T}. In 12 tested cases, the absolute discrepancies at 2,048 quadrature nodes were below the portable tolerance $2\times10^{-10}$. The accompanying machine-readable log records the measured maximum.
It also checked the two closed-form endpoint formulas at 70-digit precision in six cases, with errors below $10^{-60}$, and tested 363,810 finite-prefix signs.
The report retains coarser 256- and 512-node results to make quadrature convergence inspectable.

Representative rigorous enclosures, rounded outward, are:
\begin{center}
\small
\begin{tabular}{rrrr}
\toprule
$s$ & $v$ & $\lambda$ & Certified interval for $\Ical(s,v,\lambda)$\\
\midrule
$6$ & $3/4$ & $1/100$ & $[-0.885790290733,\,-0.885790290732]$\\
$6$ & $2$ & $1/10$ & $[-0.343926212150,\,-0.343926212149]$\\
$6$ & $24/7$ & $1$ & $[-0.043889481596,\,-0.043889481595]$\\
$20$ & $10$ & $1/20$ & $[-0.245502894584,\,-0.245502894583]$\\
\bottomrule
\end{tabular}
\end{center}
The machine-readable certificate file gives tighter 20-decimal outward enclosures for all seven examples.

\textbf{Status check.} The pinned repository entry inspected on 1 October 2026 is labeled Open and records a review date of 24 September 2026 \cite{aim}. The current arXiv record of the source paper lists version 3, dated 24 November 2025 \cite{chen}. Targeted searches for this exact variance conjecture did not identify a later resolution. Such a search cannot prove the absence of an unindexed result.

\textbf{Review status.} The supplied derivation is AI-generated. Codex read the complete argument and pinned target; separate Codex AI reviewers checked the mathematics and computational programs. No human peer review, journal acceptance, or proof-assistant verification is claimed. Review details and reproduction results accompany this note. Under the repository's evidence convention, the appropriate requested status is \emph{Solution claimed}, pending independent review \cite{contributing}. The submission includes the proof, target comparison and reproducibility records.

\appendix
\section{The exact tail bound used in the example certificates}
For rational $q,z,\lambda$ in the theorem's parameter range, set
\[
 A=1-qz,\qquad R=\sqrt{\frac{1-q^2}{1-qz}}\in(0,1).
\]
Each summand in \eqref{eq:reduction}, without its positive prefactor, is
\begin{equation*}\begin{aligned}
 w_na_n=w_nc_n\bigg[&A n z^{2n-1}-\frac q2z^{2n}\\
 &+R\left(-A n q^nz^{n-1}+\frac q2(qz)^n\right)\bigg].
\end{aligned}\end{equation*}
Terms with $n=0$ are interpreted via the derivative of $z^0$, namely zero.
Thus a finite sum has the exact form $U+RV$ with rational $U,V$.
Integer square roots bound $R$ between adjacent multiples of $10^{-50}$.

For $x=|z|>0$, put $u=q^2x^2$, $t=q^3x$, and $m=N+1$.
Because $c_n\le1$, $R<1$, and $w_n\le(1-q)q^{2n}/\lambda$, the omitted tail is at most
\begin{equation}\label{eq:tail}
 \frac{1-q}{\lambda}\left[
 \frac A{x}\bigl(F_m(u)+F_m(t)\bigr)
 +\frac q2\bigl(H_m(u)+H_m(t)\bigr)\right],
\end{equation}
where
\[
 H_m(a)=\frac{a^m}{1-a},\qquad
 F_m(a)=\frac{a^m\bigl(m-(m-1)a\bigr)}{(1-a)^2}.
\]
At $z=0$ the tail is zero once $N\ge1$.
Multiplication by an outward enclosure of
$2v z'(v)/\sqrt{1-q^2}$ produces the certified interval for $\Ical$.
All arithmetic in this procedure is integer or rational arithmetic.

\begin{thebibliography}{9}
\bibitem{aim}
S. Brunton, M. J. Colbrook, M. V. de Hoop, G. Stepaniants, A. Townsend and R. Ward.
\emph{AIM---Open Applied Problems}, Problem 560: ``Monotone variance in the Gaussian approximation to DrMMD flow.'' Repository commit \path{aa776a01d7d48a79f93251af11fde9454b0aea95}; accessed 1 October 2026.
\url{https://github.com/MColbrook/AIM/blob/aa776a01d7d48a79f93251af11fde9454b0aea95/problems/560-drmmd-gaussian-variance-monotonicity.md}.

\bibitem{chen}
Z. Chen, A. Mustafi, P. Glaser, A. Korba, A. Gretton and B. K. Sriperumbudur.
(De)-regularized Maximum Mean Discrepancy Gradient Flow.
\emph{Journal of Machine Learning Research} \textbf{26}(235) (2025), 1--77.
Appendix C, especially pp.\ 71--72 and Lemma C.2.
\url{https://www.jmlr.org/papers/v26/24-1574.html}.
Also arXiv:2409.14980v3, 24 November 2025,
\url{https://arxiv.org/abs/2409.14980v3}.

\bibitem{contributing}
AIM contributors. \emph{Contributing: Status and evidence}.
Repository commit \path{aa776a01d7d48a79f93251af11fde9454b0aea95}; accessed 1 October 2026.
\url{https://github.com/MColbrook/AIM/blob/aa776a01d7d48a79f93251af11fde9454b0aea95/CONTRIBUTING.md}.
\end{thebibliography}
\end{document}
